\chapter{Conclusão}\label{cap:conclusoeseperspectivas}

Este trabalho desenvolveu e avaliou uma infraestrutura de meta-aprendizado para ranquear pipelines de classificação com descritores estatísticos de conjuntos rotulados. A auditoria das tarefas e a reconstrução da Fase 1 permitiram revisar resultados e conclusões com maior controle experimental.

\section{Síntese dos Resultados}\label{sec:sintese-conclusao}

O acervo histórico continha 201 tarefas. Os critérios de qualidade delimitaram 193 tarefas em 165 grupos, sobre as quais foram executados 57900 ajustes dos 60 pipelines em cinco folds. Todas as etapas ajustáveis de pré-processamento usaram somente treino, e a SVM recebeu as mesmas amostras dos demais classificadores. Os avisos e as falhas foram documentados, sem substituir escores ausentes por resultados de teste.

Na avaliação principal, o Kendall ponderado médio por grupo foi $0{,}4280$ para k-NN, $0{,}4768$ para Random Forest e $0{,}4011$ para média global. A diferença Random Forest menos k-NN foi $0{,}0515$, com intervalo descritivo $[0{,}0255;0{,}0771]$. O intervalo descritivo sugere vantagem do Random Forest sobre o k-NN neste acervo. O diagnóstico encontrou 13 tarefas com Kendall ponderado não positivo e 16 com correlação indefinida.

A comparação com escores históricos usou as mesmas tarefas e partições, distinguindo a reconstrução de escores da recomputação da representação. Os valores de LOO por tarefa divulgados inicialmente não constituem uma comparação direta com a validação revisada por grupos. A melhoria do protocolo deve ser distinguida de qualquer alteração numérica de desempenho.

Na comparação pareada com a referência revisada, a mudança de Kendall ponderado foi $-0{,}0093$ para k-NN e $0{,}0059$ para Random Forest. O intervalo da mudança do k-NN ajustado abrange zero. O intervalo da mudança do Random Forest abrange zero. A elevação das médias gerais, calculadas sobre suportes diferentes, não estabelece melhoria geral da predição.

\section{Atendimento dos Objetivos}\label{sec:objetivosAtendidos}

O primeiro objetivo foi atendido pela construção e reconstrução auditável do meta-conjunto, com dados, configurações, escores e falhas rastreáveis. O segundo foi atendido pela implementação e avaliação de preditores de ranking com ajuste interno e grupos externos separados. O ganho depende da métrica e do comparador; não se conclui adequação uniforme para todos os problemas.

\section{Limitações e Perspectivas}\label{sec:limitacoes}

A representação determinística exige rótulos na consulta e usa codificações nominais que limitam a interpretação de algumas correlações. A avaliação continua retrospectiva, os grupos podem não identificar todas as versões relacionadas, os hiperparâmetros dos classificadores e do Random Forest meta foram mantidos fixos e o orçamento de MLP pode produzir avisos de convergência. Os intervalos de bootstrap são descritivos devido à sobreposição dos treinos.

A reconstrução integral da Fase 1 foi concluída nesta etapa; deixa de ser uma pendência do trabalho. As próximas etapas são ampliar o arsenal de hiperparâmetros, comparar representações sob o mesmo protocolo, avaliar orçamento e convergência e validar em dados externos coletados após o desenvolvimento. O piloto anterior de hiperparâmetros permanece uma análise exploratória separada.
